% W10: relative Chern comparison on the original labelled lattice family.
\subsection{The two Frobenius corrections by relative Chern characters}
\label{sec:relative-chern-correction}

The following argument replaces the root-by-root computation of the
two-level correction by a calculation with two short exact sequences.
It is a numerical descent theorem: its input is the original
\emph{pointwise} torsion-lattice identities, not the desired global
Chern comparison. In particular those pointwise identities still need
a proof for the actual torsion orders.

Let \(X\) be an integral, normal, projective flat model of relative
dimension \(d\) over an excellent mixed-characteristic DVR. Let
\(Z=X_{p=0}\), where the rational prime \(p\) is nonzero in the
base, and suppose that a base element \(w\) satisfies
\(w^{p-1}=up\) with \(u\) a unit. The generic fibre is denoted \(X_\eta\).
Write \(T=\mathcal O_X^n\). All lattices below carry their
specified generic identifications, and all maps between them are
the prescribed coordinate maps. We use the same
labels \(\tau\) and Frobenius permutation \(\phi\) as in the source.
The four bundles are written \(L_\tau,W_\tau,O_\tau,O_\tau^{\rm HT}\),
where \(L\) denotes the old Lie bundle and \(O\) the old
cotangent bundle, and
\[
 L_\tau\subset W_\tau\subset w^{-1}L_\tau,\qquad
 O_\tau^{\rm HT}\subset O_\tau\subset w^{-1}O_\tau^{\rm HT}.
\]
No identification of the two different corrected bundles is made.

Here is the lattice input, in coordinate notation. It is exactly the
part of \texttt{LieLatticeData.Base} and \texttt{Facts} used in this
calculation. Let \(\mathbb F\) be the finite field of the torsion
coordinates, \(M\subset\mathcal O_X^{\mathbb F^n}\) the original
function lattice, containing the constants and a power of \(p\) times
the full function lattice. Write \(\chi_\tau\) for the original
character functions with \(\chi_\tau(0)=0,\chi_\tau(1)=1\). Put
\[
 A_\tau=\{a:\ \sum_i a_i\chi_\tau(x_i)
                      \in M+p\mathcal O_X^{\mathbb F^n}\},
 \qquad H_\tau=A_\tau^\vee .
\]
Here and below all constructions on lattices are sheaf constructions;
duals use the given standard pairing. The lattice \(A_\tau\)
lies in \(T\), by evaluation at the coordinate vectors. We require
\(wT\subset A_\tau\). The original two graphs
\(\Gamma_\tau\subset L_{\tau,\eta}\oplus T_\eta\) and
\(\Delta_\tau\subset T_\eta\oplus O_{\tau,\eta}\)
have the following projections and coordinate kernels:
\[
\begin{array}{c|cc|cc}
 &\operatorname{pr}_1&\operatorname{pr}_2&
 \ker(\operatorname{pr}_1)&\ker(\operatorname{pr}_2)\\ \hline
 \Gamma_\tau &w^{-1}W_\tau&C_\tau^{\rm all}&T&W_\tau\\
 \Delta_\tau &w^{-1}T&w^{-1}O_\tau^{\rm HT}
      &O_\tau^{\rm HT}&C_\tau^{\rm all}
\end{array}
\]
where \(C_\tau^{\rm all}=\operatorname{pr}_2\Gamma_\tau\).
The two extra facts are
\[
 \operatorname{pr}_1(\Gamma_\tau\cap(L_{\tau,\eta}\oplus H_\tau))
       =w^{-1}L_\tau,\qquad
 \operatorname{pr}_2(\Delta_\tau\cap(H_\tau\oplus O_{\tau,\eta}))
       =O_\tau.                                                    \tag{RC1}
\]

For the second graph calculation let \(V_i\) be the space of functions
vanishing on the span of the first \(i\) coordinate vectors of
\(\mathbb F^n\); hence \(V_n=0\). Let \(\Theta_\tau\) be the
original graph from the augmentation lattice to the dual Lie
lattice. Its first projection is \(M\cap V_0\), its second
projection is \(L_\tau^\vee\), and its second-coordinate-only
kernel is \(p L_\tau^\vee\). Define
\[
\begin{split}
 I_{\tau i}&=\{a:\ a\chi_\tau(x_i)\in M+V_{i+1}\},\\
 K_{\tau i}&=\operatorname{pr}_2(\Theta_\tau
                         \cap(V_i\oplus L_{\tau,\eta}^\vee)),\\
 B_{\tau i}&=\{(a,b):\
       (a\chi_\tau(x_i),b)\in\Theta_\tau+(V_{i+1}\oplus0)\}.
\end{split}
\]
Indices \(i\) for the coordinates run from \(0\) to \(n-1\);
\(K_{\tau i}\) is defined also at \(i=n\).
For \(A_\tau\) use the decreasing coordinate filtration given by
\(a_j=0\) for \(j<i\), and denote the image at coordinate \(i\) by
\(D_{\tau i}\). The other three pointwise facts required are
\[
 D_{\tau i}=I_{\tau i},\qquad
 \operatorname{pr}_2 B_{\tau i}=K_{\tau i},\qquad
 \ker(\operatorname{pr}_2 B_{\phi\tau,i})
            =\langle a^p:a\in I_{\tau i}\rangle_{\mathcal O_X}.
                                                                    \tag{RC2}
\]
In the last formula the kernel is identified with its first
coordinate. This is the \emph{span of powers}, not a map
\(a\mapsto a^p\) asserted to be additive in mixed characteristic.

\paragraph{Putting the pointwise facts on one auxiliary model.}
The following elementary flattening observation allows us to start
with the source's \texttt{TestFacts} instead of assuming (RC1)--(RC2)
on the entire original model.

\begin{lemma}[Simultaneous lattice diagram]
\label{lem:relative-chern-lattice-diagram}
Suppose that the original coherent lattice family has the indicated
base projections, kernels and inclusions, with the four bundles and
\(M\) locally free. Suppose the six facts
\(wT\subset A_\tau\), (RC1) and (RC2) hold after extension to
each of the algebraically closed, complete rank-one valued fields
in the source's injective test-point hypothesis, with every
lattice interpreted as its generated lattice there.
There is a common normal projective birational modification,
an isomorphism over \(X_\eta\), on which these facts hold
as sheaf identities and all bundles, filtered pieces and graph
projections needed below are locally free. The original four
bundles on it are their pullbacks.
\end{lemma}
\begin{proof}
For any coherent fractional module of constant generic rank,
take the closure of the generic identity quotient in its relative
Grassmannian of quotients of that rank, and normalize.
The tautological quotient is the pullback modulo torsion:
the kernel is generically zero and the quotient is locally free.
Thus any finite list of coherent image lattices can be made
locally free in succession. A surjection onto a locally free
image splits locally. Consequently its kernel (and the intersection
that this kernel represents) is locally free if its domain is
locally free, and these constructions now commute with every
extension to a valued test field. More explicitly, for a fixed
linear subspace use the image in the quotient by that subspace.
For the intersection of two lattices in one space use the image
of their direct sum under subtraction, i.e.\ their sum lattice.
Once that image is locally free, the kernel is the intersection
and the short exact sequence remains exact on pullback.
The image of a further coordinate projection is treated by
one more such modification. All subsequent pullbacks preserve
the locally split sequences already obtained.

For the present finite diagram apply this procedure first to
\(M+p\mathcal O^{\mathbb F^n}\), its quotient by the range of
the linear character map, and the coordinate flag in \(A_\tau\).
This constructs \(A_\tau,D_{\tau i},H_\tau\) with their actual
generic identifications. For \(I_{\tau i}\) use the restriction
image of \(M\) on the first \(i+1\) coordinates, followed by
intersection with the line spanned by \(\chi_\tau(x_i)\).
The latter line has coordinate \(\chi_\tau(1)=1\).
For \(K_{\tau i}\) and \(B_{\tau i}\) use the restriction
images of \(\Theta_\tau\) (retaining the second coordinate for
\(B_{\tau i}\)), the corresponding kernels and the intersections
with that character line times the dual Lie space.
For (RC1) take the sums with \(w^{-1}W_\tau\oplus H_\tau\)
and \(H_\tau\oplus w^{-1}O_\tau^{\rm HT}\), respectively, and
the indicated coordinate images. Also flatten the image of
the restricted \(\Gamma_\tau\) in its second coordinate.
These are finitely many coherent images and kernels.

The original graphs have the required locally free
presentations: for \(\Gamma_\tau\) use the first projection
with kernel \(T\); for \(\Delta_\tau\) use the first projection
with kernel \(O_\tau^{\rm HT}\); and for \(\Theta_\tau\) use the
first projection with kernel \(pL_\tau^\vee\). The augmentation
lattice \(M\cap V_0\) is a direct summand since \(1\in M\).
The second projection of \(\Delta_\tau\) is onto the locally
free \(w^{-1}O_\tau^{\rm HT}\), so its first-coordinate kernel
is also locally free. This is \(C_\tau^{\rm all}\), and hence
the second projection sequence of \(\Gamma_\tau\) is locally
split as well. Thus the base identities survive all pullbacks.
Over \(X_\eta\) the graphs are the coordinate direct products
and the function lattice is full; all the above images have
constant rank there. The modifications are therefore vertical.

On the final normal model take any height-one point over \(p\).
Its local ring is a DVR. Extend its fraction field, via completion
and completed algebraic closure, to an algebraically closed
complete rank-one valued field. The valuation ring is faithfully
flat over that DVR: it is torsion free and the local maximal
ideal has proper extension. All the indicated images, intersections
and duals commute with this extension by the split sequences.
The resulting test field contains the original function field
injectively and has positive valuation of \(p\); the original
\texttt{TestFacts} applies to the original chart and its
generated lattices. Faithful flatness yields the required equalities
or inclusions over the height-one local ring itself.
For the last equality in (RC2), a locally free fractional
line is principal on each trivialization; its span of \(p\)-th
powers is its ordinary product ideal of order \(p\), also
locally free. Thus both sides of each required comparison
are locally free lattices. On a normal scheme such lattices
are the intersections of their height-one localizations
inside their rational vector space. All facts follow globally,
including at points of codimension greater than one.
\end{proof}

The auxiliary modification in the lemma is used only for
intersection equalities. No uniform degree, height or number of
charts of it is needed. In particular it is not a replacement
for the prescribed tower or its sections.

\paragraph{Relative numerical classes.}
For two bundles \(E,F\) with the same generic identification,
write \([E:F]\) for the relative \(K\)-class supported on
\(Z\), with sign \(E-F\). If \(F\subset E\) it is represented
by the two-term resolution of \(E/F\). In general one can
use a common sublattice \(p^a E\subset E\cap F\) and subtract
the two quotient classes; additivity makes this independent
of \(a\). All exact sequences used below respect their generic
identifications. We use the localized Chern character of
perfect complexes (Fulton, \emph{Intersection Theory},
second edition, \S18.1, Theorem 18.1, Proposition 18.1 and
Corollary 18.1.2). Its additivity and projection formula apply
to the proper models over the DVR, including singular
special fibres. For an original complementary monomial \(Q\)
of degree \(d+1-k\) in line classes, write
\[
 c_k(E:F;Q)=k!\deg_Z
       \bigl(Q\,\operatorname{ch}^{Z}_k([E:F])\cap[X]\bigr).
\]
The degree is taken on the supporting fibre, \emph{before}
pushing forward to \(X\). Forgetting the support would make
the principal scalar divisors vanish incorrectly.

For a pair of lines \(E,F\) the expression is the vertical
divisor of their generic ratio intersected with
\(\sum_{j=0}^{k-1}c_1(E)^jc_1(F)^{k-1-j}Q\).
This follows by applying the localized character to the
two-term line inclusion (or subtracting two inclusions
through a common smaller line). Hence, by additivity along
matched flags, \(c_k(E:F;Q)\) is precisely the source's
mixed-difference power sum \texttt{ck}, including the sign
\(x-x'\). Two consequences used here are
\[
\begin{split}
 c_k(J^{\otimes p}:\mathcal O;Q)
   &=p^k c_k(J:\mathcal O;Q) \quad(J\text{ a fractional line}),\\
 c_k(w^{-1}E:E;Q)
   &= k\,\deg_Z\bigl(Q P_{k-1}(E)\cap[\operatorname{div}(w)]\bigr),
                                                               \tag{RC3}
\end{split}
\]
where \(P_j(E)=j!\operatorname{ch}_j(E)\).
For the second identity tensor the Cartier divisor resolution
by \(E\): its normal line \(\mathcal O(\operatorname{div}(w))\)
is trivial, so there are no higher self-intersection terms.
For the first identity use the line formula and
\(c_1(J^{\otimes p})=p c_1(J)\), with the ratio divisor
also multiplied by \(p\). In particular the normalized
power-sum factor \(k\) in the second line is necessary.

\begin{theorem}[Two relative Chern identities and the finite orbit formula]
\label{thm:relative-chern-two-identities}
Under the base and test hypotheses of
Lemma~\ref{lem:relative-chern-lattice-diagram}, set, for the original
complementary monomial \(Q\),
\[
 h_\tau=c_k(H_\tau:T;Q),\qquad
 \ell_\tau=k\deg_Z\bigl(QP_{k-1}(L_\tau)
                              \cap[\operatorname{div}(p)]\bigr).
\]
For every \(1\le k\le d+1\) the original two-level numerical
relation is
\[
\begin{split}
 c_k(L_\tau:W_\tau;Q)+c_k(O_\tau:O_\tau^{\rm HT};Q)
     &=h_\tau-\ell_\tau/(p-1),\\
 p^kh_\tau&=h_{\phi\tau}+\ell_{\phi\tau}.                 \tag{RC4}
\end{split}
\]
If the orbit of \(\tau\) under \(\phi\) has length \(m\), then
\[
 h_\tau=\frac{\sum_{a=1}^{m}p^{k(m-a)}
                                  \ell_{\phi^a\tau}}{p^{km}-1}.
                                                               \tag{RC5}
\]
These are the exact \texttt{Relations} and
\texttt{LabelChernComparison}, not estimates or asymptotic
formulas. They descend to the original model and hold with
any original homogeneous complementary polynomial by linearity.
\end{theorem}
\begin{proof}
Work first on the common model. Let
\[
 C_\tau=\operatorname{pr}_2\bigl(\Gamma_\tau
                  \cap(L_{\tau,\eta}\oplus H_\tau)\bigr).
\]
The restricted graph has first projection \(w^{-1}L_\tau\),
first-coordinate-only kernel \(W_\tau\), and
second-coordinate-only kernel \(T\). Taking the quotient by
\(W_\tau\oplus T\) in two ways gives
\[
 C_\tau/T\simeq w^{-1}L_\tau/W_\tau.
\]
The restricted second graph has first projection \(H_\tau\)
because \(H_\tau\subset w^{-1}T\); its second projection
is \(O_\tau\). Its two coordinate kernels are
\(C_\tau\) and \(O_\tau^{\rm HT}\). Taking the quotient
by their direct sum gives
\(H_\tau/C_\tau\simeq O_\tau/O_\tau^{\rm HT}\).
The three quotients are supported on \(Z\). Additivity in
relative \(K\)-theory therefore proves the exact identity
\[
 [H_\tau:T]=[L_\tau:W_\tau]+[O_\tau:O_\tau^{\rm HT}]
                         +[w^{-1}L_\tau:L_\tau].            \tag{RC6}
\]
Apply (RC3) and \((p-1)\operatorname{div}(w)=\operatorname{div}(p)\).
This gives the first equality of (RC4).

For the other identity, the first projection of \(B_{\tau i}\)
is \(I_{\tau i}\): a witness in \(M\) restricting to
\(a\chi_\tau(x_i)\) vanishes at zero, so has a preimage in
\(\Theta_\tau\). The second-coordinate-only kernel is
\(K_{\tau,i+1}\). By (RC2) the second projection is
\(K_{\tau i}\) and the first-coordinate-only kernel is
\(I_{\phi^{-1}\tau,i}^{\otimes p}\).
Taking the quotient of the graph by the two kernels gives
\[
 I_{\tau i}/I_{\phi^{-1}\tau,i}^{\otimes p}
                       \simeq K_{\tau i}/K_{\tau,i+1}.
\]
All these are quotients of locally free lattices, and the
isomorphism is induced by the original graph. Equivalently
one can compare its two locally split projection sequences.
Dualizing those sequences, with the generic coordinate
identifications unchanged, gives
\[
 [I_{\tau i}^{\vee}:(I_{\phi^{-1}\tau,i}^{\vee})^{\otimes p}]
                   =[K_{\tau i}^{\vee}:K_{\tau,i+1}^{\vee}].
                                                               \tag{RC7}
\]
This sign is important: both dual comparisons reverse the
respective inclusions. The coordinate filtration of \(A_\tau\)
has graded lines \(D_{\tau i}=I_{\tau i}\); the dual filtration
gives
\([H_\tau:T]=\sum_i[I_{\tau i}^\vee:\mathcal O]\).
Also \(K_{\tau0}=L_\tau^\vee\), \(K_{\tau n}=pL_\tau^\vee\),
by the two projection conditions of \(\Theta_\tau\).
Sum (RC7), apply the first identity of (RC3), and telescope
the right side. The result is
\[
 h_\tau-p^kh_{\phi^{-1}\tau}
   =c_k(L_\tau:p^{-1}L_\tau;Q)=-\ell_\tau,
\]
the second equality of (RC4) after relabelling.

Iterate this equality exactly \(m\) times; since \(p^{km}>1\),
division gives (RC5). No infinite tower of Chern identities
or convergence of model corrections is involved.
All auxiliary pullbacks of the original class comparisons
have unchanged numerical values by the localized projection
formula, since the birational map pushes the fundamental
cycle to the original one with degree one.
The same formula applies to \(\operatorname{div}(p)\)
and the old bundle roots. Hence (RC4)--(RC5) prove the
numerical identities on the original model.
\end{proof}

For clarity, the coefficient of \(\ell_{\tau'}\) in the
first line of (RC4) is zero outside the orbit and within
the orbit is exactly
\[
 \frac{p^{k(m-a)}}{p^{km}-1}
             -\frac{\mathbf1_{\tau'=\tau}}{p-1},
 \qquad \tau'=\phi^a\tau,\quad 1\le a\le m .
\]
This agrees with the source's \texttt{co}: group its geometric
series into complete periods. In particular the absolute row sum
is
\[
 \frac1{p-1}+\frac1{p^k-1}-\frac2{p^{km}-1}.
\]
Summing over a complete orbit gives the simpler identity
\[
 \sum_{\tau\ {\rm in\ orbit}}
  (c_k(L_\tau:W_\tau;Q)+c_k(O_\tau:O_\tau^{\rm HT};Q))
 =\left(\frac1{p^k-1}-\frac1{p-1}\right)
                      \sum_{\tau\ {\rm in\ orbit}}\ell_\tau .
\]
In particular the aggregate first correction is zero on
each complete orbit for \(k=1\). This cancellation does not
say that a single labelled first correction is zero.

With the original dual and sign identifications of the
corrected Hodge flags, the old Lie power sums in the labels
\(\tau,\bar\tau\) are exactly the two entries of \texttt{liePs},
and the left side of (RC4) is the original full flag power
sum difference for the corresponding label. Thus the formula
has precisely the coefficients of \texttt{FlagChernComparison}.
This final identification must use the actual labelled
bundle pairs at each correction prime, not merely bundles
with equal ranks. The previously established trivial
corrected-flag theorem handles the separate fibre-vanishing
condition. This calculation leaves the actual test-point
torsion-order facts, the global two-level model linkage,
the weighted tower mean and the section-rank losses to be
completed; it does not establish W10 or the final RH budget.
